\subsubsection{Subproceso 2: Selección y Registro de Cursos}



\begin{itemize}
    \item \textbf{Estudiante Regular:} Actor principal y único usuario con capacidad de decisión sobre la selección de asignaturas. Es responsable de conformar su horario dentro del marco normativo, respetar el límite crediticio permitido y confirmar el envío de la propuesta dentro de su turno habilitado.

    \item \textbf{Intranet DIRCE / Sistema SIGA-ORCE:} Actor tecnológico que opera como motor de validación y control de concurrencia. Gestiona la disponibilidad de cupos en tiempo real, aplica los filtros de prerrequisitos y cruces de horario de manera automática, y persiste la Pre-Ficha una vez confirmada la selección.

    \item \textbf{Comisión de Matrícula:} Órgano colegiado de la Facultad encargado de parametrizar la oferta académica semestral (cursos, secciones, horarios, docentes y cupos) antes del inicio del proceso. Su labor en esta etapa es de carácter previo y de soporte; no interviene directamente en la selección del estudiante, pero sus decisiones de programación condicionan la oferta visible en el portal.

    \item \textbf{Departamentos Académicos:} Unidades responsables de la asignación de docentes y del armado de los horarios por sección, los cuales son ingresados al SIGA-ORCE como insumo de la oferta que el estudiante visualiza en esta etapa.
\end{itemize}


